% ============================================
% USTC PhD Exam Prep - Main Document
% 中科大物理博士考试备考资料
% ============================================
% ============================================
% USTC PhD Exam Prep - Preamble
% 中科大物理博士考试备考资料导言区
% ============================================

\documentclass[10pt,a4paper]{article}

% ---------- 基础包 ----------
\usepackage[margin=1.2cm, top=1cm, bottom=1cm]{geometry}
\usepackage{ctex}  % 中文支持
\usepackage{multicol}  % 双栏/三栏
\usepackage{enumitem}  % 列表控制
\usepackage{hyperref}  % 超链接

% ---------- 数学包 ----------
\usepackage{amsmath,amssymb,amsthm}
\usepackage{physics}
\usepackage{mathrsfs}
\usepackage{bm}  % 粗体数学符号

% ---------- 图形与颜色 ----------
\usepackage{xcolor}
\usepackage{tcolorbox}
\tcbuselibrary{skins,breakable}
\usepackage{tikz}
\usetikzlibrary{shapes,arrows,positioning,calc}

% ---------- 页面设置 ----------
\pagestyle{empty}  % 无页眉页脚，最大化空间
\setlength{\parindent}{0pt}  % 无缩进
\setlength{\parskip}{2pt}    % 紧凑段落
\setlength{\headheight}{20pt}  % 修复 fancyhdr 警告

% ---------- 双栏设置 ----------
\setlength{\columnseprule}{0.4pt}  % 栏间分隔线
\def\columnseprulecolor{\color{gray!50}}

% ============================================
% 颜色定义 - Beamer 风格
% ============================================

% 主题色
\definecolor{ustcblue}{RGB}{0,51,153}      % 科大蓝
\definecolor{alertred}{RGB}{220,38,38}     % 警示红
\definecolor{successgreen}{RGB}{34,139,34} % 成功绿
\definecolor{warnorange}{RGB}{255,140,0}   % 警告橙
\definecolor{royalpurple}{RGB}{120,81,169} % 皇家紫

% 柔和背景色
\definecolor{lightblue}{RGB}{230,240,255}
\definecolor{lightred}{RGB}{255,230,230}
\definecolor{lightgreen}{RGB}{230,255,230}
\definecolor{lightorange}{RGB}{255,245,230}
\definecolor{lightpurple}{RGB}{245,230,255}
\definecolor{lightyellow}{RGB}{255,255,200}

% ============================================
% tcolorbox 样式定义
% ============================================

% 1. 核心定理/模型框（蓝色主题）
\newtcolorbox{modelbox}[1][]{
    enhanced,
    colback=lightblue,
    colframe=ustcblue,
    fonttitle=\bfseries\sffamily\small,
    title=#1,
    boxrule=1pt,
    arc=3pt,
    left=4pt,right=4pt,top=3pt,bottom=3pt,
    drop shadow={opacity=0.15},
    breakable
}

% 2. 解题技巧/快速识别框（红色主题）
\newtcolorbox{tipbox}[1][]{
    enhanced,
    colback=lightred,
    colframe=alertred,
    fonttitle=\bfseries\sffamily\small,
    title=#1,
    boxrule=1pt,
    arc=3pt,
    left=4pt,right=4pt,top=3pt,bottom=3pt,
    drop shadow={opacity=0.15},
    breakable
}

% 3. 重要公式/二级结论框（绿色主题）
\newtcolorbox{formulabox}[1][]{
    enhanced,
    colback=lightgreen,
    colframe=successgreen,
    fonttitle=\bfseries\sffamily\small,
    title=#1,
    boxrule=1pt,
    arc=3pt,
    left=4pt,right=4pt,top=3pt,bottom=3pt,
    drop shadow={opacity=0.15},
    breakable
}

% 4. 警示/易错点框（橙色主题）
\newtcolorbox{warnbox}[1][]{
    enhanced,
    colback=lightorange,
    colframe=warnorange,
    fonttitle=\bfseries\sffamily\small,
    title=#1,
    boxrule=1pt,
    arc=3pt,
    left=4pt,right=4pt,top=3pt,bottom=3pt,
    drop shadow={opacity=0.15},
    breakable
}

% 5. 物理图像/概念理解框（紫色主题）
\newtcolorbox{insightbox}[1][]{
    enhanced,
    colback=lightpurple,
    colframe=royalpurple,
    fonttitle=\bfseries\sffamily\small,
    title=#1,
    boxrule=1pt,
    arc=3pt,
    left=4pt,right=4pt,top=3pt,bottom=3pt,
    drop shadow={opacity=0.15},
    breakable
}

% 6. 速查/备忘框（黄色主题）
\newtcolorbox{cheatbox}[1][]{
    enhanced,
    colback=lightyellow,
    colframe=warnorange!80!black,
    fonttitle=\bfseries\sffamily\small,
    title=#1,
    boxrule=1pt,
    arc=3pt,
    left=3pt,right=3pt,top=2pt,bottom=2pt,
    breakable
}

% 7. 附录推导框（灰色主题）
\newtcolorbox{derivationbox}[1][]{
    enhanced,
    colback=gray!10,
    colframe=gray!70!black,
    fonttitle=\bfseries\sffamily\small,
    title=#1,
    boxrule=0.8pt,
    arc=2pt,
    left=3pt,right=3pt,top=2pt,bottom=2pt,
    breakable
}

% ============================================
% 快捷命令
% ============================================

% 英文术语备注
\newcommand{\en}[1]{\textcolor{ustcblue}{\textit{#1}}}

% 重要强调
\newcommand{\key}[1]{\textcolor{alertred}{\textbf{#1}}}

% 公式中的文字说明
\newcommand{\texteq}[1]{\quad\text{\small #1}}

% 快速引用
\newcommand{\quickref}[2]{\fbox{\parbox{0.9\linewidth}{\small\textbf{#1:} #2}}}

% 考试提醒标记
\newcommand{\examnote}[1]{\textcolor{warnorange}{$\triangleright$ \textbf{#1}}}

% ============================================
% 列表样式优化
% ============================================
\setlist[itemize]{
    nosep,
    leftmargin=12pt,
    label={\textcolor{ustcblue}{$\bullet$}}
}

\setlist[enumerate]{
    nosep,
    leftmargin=15pt
}

% ============================================
% 章节标题样式
% ============================================
\usepackage{titlesec}

\titleformat{\section}
    {\Large\bfseries\sffamily\color{ustcblue}}
    {\thesection}{0.5em}{}
    [\vspace{-0.5em}\textcolor{ustcblue}{\rule{\linewidth}{1pt}}]

\titleformat{\subsection}
    {\large\bfseries\sffamily\color{alertred}}
    {\thesubsection}{0.5em}{}

\titleformat{\subsubsection}
    {\normalsize\bfseries\sffamily\color{successgreen}}
    {\thesubsubsection}{0.5em}{}

\titlespacing*{\section}{0pt}{10pt}{5pt}
\titlespacing*{\subsection}{0pt}{8pt}{4pt}
\titlespacing*{\subsubsection}{0pt}{6pt}{3pt}

% ============================================
% 页眉设置（可选）
% ============================================
\usepackage{fancyhdr}
\pagestyle{fancy}
\fancyhf{}
\fancyhead[L]{\small\textcolor{gray}{USTC PhD Exam Prep}}
\fancyhead[R]{\small\textcolor{gray}{\leftmark}}
\renewcommand{\headrulewidth}{0.4pt}
\renewcommand{\headrule}{\hbox to\headwidth{\color{gray}\leaders\hrule height \headrulewidth\hfill}}


\begin{document}

% ---------- 封面 ----------
\begin{center}
{\Huge\bfseries\sffamily\color{ustcblue} 中科大物理博士考试备考资料}\\[8pt]
{\Large USTC PhD Physics Exam Preparation}\\[12pt]
{\large 整理者：Qi Meng}\\[4pt]
{\normalsize 更新日期：\today}
\end{center}

\vspace{8pt}
\textcolor{ustcblue}{\rule{\linewidth}{2pt}}
\vspace{8pt}

% ============================================
% 核心页：压缩算符与标度变换
% ============================================
\section{量子力学：压缩算符与标度变换}

\begin{multicols}{2}

\begin{modelbox}[题目：压缩算符的标度变换]
定义压缩算符（\en{squeezing operator}）：
\[
S(\xi) = \exp\left\{ \frac{1}{2}\left( \xi^* a^2 - \xi (a^\dagger)^2 \right) \right\}
\]
其中 $a$ 和 $a^\dagger$ 是谐振子的升降算符。若取 $\xi = r$ 为实数，证明 $S(r)$ 为标度变换，即证明 $S^\dagger(r) X S(r)$ 正比于 $X$，$S^\dagger(r) P S(r)$ 正比于 $P$；证明两者的比例系数互为倒数，并求出具体的比例系数。
\end{modelbox}

\begin{tipbox}[考点快速分析]
\begin{itemize}
    \item \textbf{算符代数}：指数算符酉变换 $e^{-G} A e^{G}$ 的微分方程处理
    \item \textbf{压缩态}：相空间中的几何形变（一个方向压缩、另一个方向膨胀）
    \item \textbf{Bogoliubov 变换}：玻色子配对哈密顿量的对角化基础
\end{itemize}
\end{tipbox}

\begin{derivationbox}[标准解题过程]
\textbf{Step 1: 算符关系构建}

在自然单位 $\hbar = m = \omega = 1$ 下：
\[
X = \frac{1}{\sqrt{2}}(a + a^\dagger), \quad P = \frac{-i}{\sqrt{2}}(a - a^\dagger)
\]

\textbf{Step 2: 转化为微分方程}

令 $a(r) = S^\dagger(r) a S(r)$。由于 $S(r) = e^{rG}$ 且 $G = \frac{1}{2}(a^2 - (a^\dagger)^2)$，有
\[
\frac{da(r)}{dr} = [a(r), G]
\]
计算对易子：
\[
[a, G] = \frac{1}{2}\left[ a, a^2 - (a^\dagger)^2 \right] = -a^\dagger
\]
建立耦合方程组：
\[
\begin{cases}
\displaystyle \frac{da(r)}{dr} = -a^\dagger(r) \\[6pt]
\displaystyle \frac{da^\dagger(r)}{dr} = -a(r)
\end{cases}
\]

\textbf{Step 3: 求解并代入}

双曲函数解（初始条件 $a(0)=a$）：
\[
a(r) = a\cosh r - a^\dagger \sinh r
\]
\[
a^\dagger(r) = a^\dagger \cosh r - a \sinh r
\]

对于 $X$：
\begin{align*}
S^\dagger X S &= \frac{1}{\sqrt{2}}(a(r) + a^\dagger(r)) \\
&= \frac{1}{\sqrt{2}}(a + a^\dagger)(\cosh r - \sinh r) = e^{-r} X
\end{align*}

对于 $P$：
\begin{align*}
S^\dagger P S &= \frac{-i}{\sqrt{2}}(a(r) - a^\dagger(r)) \\
&= \frac{-i}{\sqrt{2}}(a - a^\dagger)(\cosh r + \sinh r) = e^{r} P
\end{align*}
\end{derivationbox}

\begin{formulabox}[核心结论]
\[
S^\dagger(r) X S(r) = e^{-r} X, \qquad S^\dagger(r) P S(r) = e^{r} P
\]
比例系数 $e^{-r}$ 与 $e^{r}$ 互为倒数，乘积为 $1$。
\end{formulabox}

\begin{errorbox}[colback=white]{思路短路原因与解决策略}
\textbf{短路原因：}
\begin{itemize}
    \item \textbf{公式焦虑}：试图直接展开 $e^A B e^{-A}$ 的无限级数（BCH 公式），导致中间项计算极其混乱
    \item \textbf{系数混淆}：在 $X \sim a+a^\dagger$ 中纠结具体的 $1/\sqrt{2}$ 或 $i$ 的正负号
    \item \textbf{静态思维}：将 $r$ 看作死板的参数，而非动态演化的``时间''
\end{itemize}

\textbf{解决策略：}
\begin{itemize}
    \item \textbf{微分化思维}：看到 $e^{r(\dots)}$ 作用于算符，第一反应永远是求导建方程。这比硬算对易子快且准
    \item \textbf{线性保持原则}：记住压缩变换是线性的，$X$ 变换后一定还是 $X$ 的线性组合
    \item \textbf{结果逆推}：既然是``压缩''，必然一个方向缩 $e^{-r}$，另一个扩 $e^{r}$
\end{itemize}
\end{errorbox}

\begin{insightbox}[举一反三与物理直觉]
\textbf{物理直觉——相空间``揉面团''：}

把真空态在相空间中想象成一个圆（相干态的准概率分布）。压缩算符 $S(\xi)$ 所生成的变换是一种保辛结构（保面积）的仿射变换：它将圆挤压成椭圆，但满足最小不确定关系 $\Delta X \Delta P = \hbar/2$。当 $\xi = r$ 为实数时，椭圆的长轴和短轴恰好与 $X$ 轴、$P$ 轴对齐，$X$ 方向被压缩 $e^{-r}$ 而 $P$ 方向被拉伸 $e^{r}$。其来源是哈密顿量中的二次项 $a^2 - (a^\dagger)^2$ 打破了谐振子基态的旋转对称性，导致粒子对的产生与湮灭关联，从而在一个方向上实现了不确定度的塌缩。

\textbf{拓展延伸：}

\textit{(1) 复数参数} $\xi = r e^{i\theta}$：
此时 $S^\dagger(\xi) a S(\xi) = a \cosh r - a^\dagger e^{i\theta} \sinh r$。定义旋转后的坐标轴
\[
X_{\theta/2} = X \cos\frac{\theta}{2} + P \sin\frac{\theta}{2}, \quad
P_{\theta/2} = -X \sin\frac{\theta}{2} + P \cos\frac{\theta}{2}
\]
则有
\[
S^\dagger(\xi) X_{\theta/2} S(\xi) = e^{-r} X_{\theta/2}, \quad
S^\dagger(\xi) P_{\theta/2} S(\xi) = e^{r} P_{\theta/2}
\]
标度变换的方向不再是坐标轴 $X, P$，而是相空间中旋转了 $\theta/2$ 角度的倾斜轴。

\textit{(2) BdG 变换联系：}玻色子配对哈密顿量
\[
H = \omega a^\dagger a + \frac{\Delta}{2}\left( a^2 + (a^\dagger)^2 \right)
\]
的对角化，本质上就是这个压缩变换的玻色子版本。关键区别在于变换系数的约束来源不同：
\begin{itemize}[leftmargin=12pt, nosep]
    \item \textbf{玻色子}：$[b, b^\dagger] = 1$ 要求 $u^2 - v^2 = 1$，用双曲函数 $(\cosh r, \sinh r)$
    \item \textbf{费米子}：$\{c, c^\dagger\} = 1$ 要求 $u^2 + v^2 = 1$，用三角函数 $(\cos\theta, \sin\theta)$
\end{itemize}
\end{insightbox}

\begin{thinkbox}[解题思路模板（算符酉变换）]
\begin{enumerate}
    \item \textbf{识别生成元}：写出 $G$ 并确认参数对应的动力学``时间''
    \item \textbf{建立微分方程}：$\displaystyle \frac{dA(\lambda)}{d\lambda} = [A(\lambda), G]$
    \item \textbf{利用线性性}：猜测解的形式（通常为双曲/三角函数线性组合）
    \item \textbf{回代验证}：将解代回原始算符定义，检查系数与对易关系
    \item \textbf{极限检查}：$\lambda \to 0$ 时是否恢复恒等变换？
\end{enumerate}
\end{thinkbox}

\end{multicols}

% ============================================
% 补充说明页
% ============================================
\newpage
\section{补充说明：压缩算符问题的完整细节}

\subsection{题型反射弧标签}

\begin{center}
\begin{tabular}{|c|l|}
\hline
\textbf{层级} & \textbf{识别标签} \\
\hline
一级（领域） & 量子力学，谐振子代数 \\
二级（方法） & 算符酉变换，指数映射求导法 \\
三级（关键词） & squeezing operator, $S^\dagger a S$, Bogoliubov transformation \\
\hline
\end{tabular}
\end{center}

\examnote{看到 $S^\dagger(r) a S(r)$ 或 $e^{-G} A e^G$，立即启动``参数求导 $+$ 微分方程''套路，不要硬算 BCH。}

\subsection{自然单位与非自然单位的对照}

在谐振子问题中，严格的关系式为（\en{不取自然单位}）：
\[
a = \sqrt{\frac{m\omega}{2\hbar}} \left( X + \frac{i}{m\omega} P \right), \qquad
a^\dagger = \sqrt{\frac{m\omega}{2\hbar}} \left( X - \frac{i}{m\omega} P \right)
\]
反解得：
\[
X = \sqrt{\frac{\hbar}{2m\omega}} (a + a^\dagger), \qquad
P = i\sqrt{\frac{\hbar m\omega}{2}} (a^\dagger - a)
\]
当且仅当取自然单位 $\hbar = m = \omega = 1$ 时，才简化为 $X = (a + a^\dagger)/\sqrt{2}$，$P = -i(a - a^\dagger)/\sqrt{2}$。在考试中如果未声明自然单位，\key{必须保留 $\sqrt{\hbar/(2m\omega)}$ 等因子}，最终标度系数不变（仍为 $e^{\pm r}$），因为系数在标度变换中会被约去。

\subsection{微分方程的来源与推导}

设 $S(r) = e^{rG}$，其中 $G = \frac{1}{2}(a^2 - (a^\dagger)^2)$。定义海森堡型演化算符 $a(r) \equiv S^\dagger(r) a S(r) = e^{-rG} a e^{rG}$。对参数 $r$ 求导：
\begin{align*}
\frac{da(r)}{dr} &= \frac{d}{dr}\left( e^{-rG} a e^{rG} \right) \\
&= (-G) e^{-rG} a e^{rG} + e^{-rG} a e^{rG} G \\
&= e^{-rG} a e^{rG} G - G e^{-rG} a e^{rG} \\
&= [a(r), G]
\end{align*}
这正是 Heisenberg 方程的类比。\key{关键技巧}：$G$ 不显含 $r$（除前面的系数 $r$ 外），因此可以直接把 $r$ 当作``时间''来求导。

\subsection{微分方程的详细解法}

由对易关系 $[a, a^\dagger] = 1$ 得：
\[
[a, G] = \frac{1}{2}\left[ a, a^2 - (a^\dagger)^2 \right] = -a^\dagger, \qquad
[a^\dagger, G] = -a
\]
于是耦合方程组为：
\[
\frac{da(r)}{dr} = -a^\dagger(r), \qquad \frac{da^\dagger(r)}{dr} = -a(r)
\]
对第一式再求导一次：
\[
\frac{d^2 a(r)}{dr^2} = -\frac{da^\dagger(r)}{dr} = a(r)
\]
即 $d^2 a/dr^2 - a = 0$。特征方程 $\lambda^2 - 1 = 0$，通解为 $a(r) = A \cosh r + B \sinh r$。由初始条件：
\begin{itemize}
    \item $r = 0$ 时 $a(0) = a$，故 $A = a$
    \item $r = 0$ 时 $\displaystyle \left.\frac{da}{dr}\right|_{r=0} = -a^\dagger(0) = -a^\dagger$，故 $B = -a^\dagger$
\end{itemize}
因此：
\[
a(r) = a \cosh r - a^\dagger \sinh r, \qquad a^\dagger(r) = a^\dagger \cosh r - a \sinh r
\]

\subsection{考察的完整知识点}

\begin{enumerate}
    \item \textbf{升降算符与坐标动量的关系}：$a, a^\dagger$ 与 $X, P$ 的线性映射（含非自然单位形式）
    \item \textbf{基本对易关系}：$[a, a^\dagger] = 1$，$[a, a] = [a^\dagger, a^\dagger] = 0$
    \item \textbf{高阶对易子}：$[a, (a^\dagger)^2] = 2a^\dagger$，$[a^2, a^\dagger] = 2a$
    \item \textbf{指数算符的求导公式}：$\displaystyle \frac{d}{dr} e^{-rG} A e^{rG} = [e^{-rG} A e^{rG}, G]$
    \item \textbf{双曲函数恒等式}：$\cosh r - \sinh r = e^{-r}$，$\cosh r + \sinh r = e^{r}$
    \item \textbf{物理图像}：不确定度在 $X$ 方向被压缩 $e^{-r}$ 倍，在 $P$ 方向被膨胀 $e^{r}$ 倍，满足 $\Delta X \Delta P = \hbar/2$（最小不确定态）
\end{enumerate}

\subsection{BdG 变换与压缩变换的数学联系}

考虑玻色子配对哈密顿量：
\[
H = \omega a^\dagger a + \frac{\Delta}{2}\left( a^2 + (a^\dagger)^2 \right)
\]
这正是\en{压缩态}的生成哈密顿量。为了对角化 $H$，引入幺正变换 $S(r) = \exp\left\{ \frac{r}{2}(a^2 - (a^\dagger)^2) \right\}$，参数 $r$ 满足：
\[
\tanh(2r) = \frac{\Delta}{\omega}
\]
经过变换后，$b \equiv S^\dagger(r) a S(r) = a \cosh r - a^\dagger \sinh r$，哈密顿量对角化为：
\[
S^\dagger(r) H S(r) = E b^\dagger b + \text{const}, \qquad E = \sqrt{\omega^2 - \Delta^2}
\]
这与费米子版本的 BdG 变换（如超导 BCS 理论中的粒子-空穴变换 $c_k = u_k \gamma_k + v_k \gamma_{-k}^\dagger$）\key{在数学结构上完全平行}：两者都是通过线性混合产生/湮灭算符来消除哈密顿量中的非对角配对项。关键区别在于变换系数的约束来源不同：
\begin{itemize}
    \item \textbf{玻色子}：为了保持对易关系 $[b, b^\dagger] = 1$，由 $b = u a + v a^\dagger$ 可得 $u^2 - v^2 = 1$，因此用双曲函数 $(\cosh r, \sinh r)$ 参数化；
    \item \textbf{费米子}：为了保持反对易关系 $\{c, c^\dagger\} = 1$，由 $c = u \gamma + v \gamma^\dagger$ 可得 $u^2 + v^2 = 1$，因此用三角函数 $(\cos\theta, \sin\theta)$ 参数化。
\end{itemize}

\subsection{拓展延伸}

\begin{itemize}
    \item \textbf{复数参数} $\xi = r e^{i\theta}$：此时 $S^\dagger(\xi) a S(\xi) = a \cosh r - a^\dagger e^{i\theta} \sinh r$。定义旋转后的坐标算符 $X_{\theta/2} = X \cos(\theta/2) + P \sin(\theta/2)$ 和 $P_{\theta/2} = -X \sin(\theta/2) + P \cos(\theta/2)$，则标度变换沿倾斜轴进行：$S^\dagger(\xi) X_{\theta/2} S(\xi) = e^{-r} X_{\theta/2}$，$S^\dagger(\xi) P_{\theta/2} S(\xi) = e^{r} P_{\theta/2}$。此时 $S^\dagger(\xi) X S(\xi)$ 和 $S^\dagger(\xi) P S(\xi)$ 一般不再是单独的 $X$ 或 $P$，而是两者的线性组合。
\end{itemize}

% ============================================
% 核心页：线性势场与相互作用绘景
% ============================================
\newpage
\section{量子力学：线性势场与相互作用绘景}

\begin{multicols}{2}

\begin{modelbox}[题目：线性势中的粒子演化]
考虑在线性势中的粒子，哈密顿量为：
\[
H = \frac{P^2}{2m} - \gamma X \quad (\gamma \in \mathbb{R})
\]
写出相互作用绘景下态矢量 $|\tilde{\psi}(t)\rangle$ 和演化算符 $\tilde{U}(t)$ 的微分方程。证明该模型存在如下形式的解，并确定相因子 $\alpha(t)$：
\[
\tilde{U}(t) = \exp\left( -\frac{i\alpha(t)}{\hbar} \right) \exp\left( \frac{i\gamma P t^2}{2m\hbar} \right) \exp\left( \frac{i\gamma X t}{\hbar} \right)
\]
最后给出薛定谔绘景中的时间演化算符 $U(t)$。
\end{modelbox}

\begin{tipbox}[考点快速分析]
\begin{itemize}
    \item \textbf{绘景拆分}：$H_0 = \frac{P^2}{2m}$（自由项），$V = -\gamma X$（作用项）
    \item \textbf{相互作用势的演化}：$V_I(t) = -\gamma (X + \frac{P}{m}t)$
    \item \textbf{对易子闭合性}：$[X, P] = i\hbar$ 为常数，保证了算符拆解的严格有限性
    \item \textbf{考场策略}：题目已给出解的形式，优先使用\key{代入验证法}而非从头推导
\end{itemize}
\end{tipbox}

\begin{derivationbox}[标准解题过程]
\textbf{Step 1: 相互作用绘景的建立}

取 $H_0 = \frac{P^2}{2m}$，则相互作用绘景态矢量定义为：
\[
|\tilde{\psi}(t)\rangle = e^{iH_0 t/\hbar} |\psi(t)\rangle
\]
其满足的运动方程为：
\[
i\hbar \frac{d}{dt}|\tilde{\psi}(t)\rangle = V_I(t) |\tilde{\psi}(t)\rangle, \quad V_I(t) = e^{iH_0 t/\hbar} V e^{-iH_0 t/\hbar}
\]
演化算符 $\tilde{U}(t)$ 满足：
\[
i\hbar \frac{d\tilde{U}}{dt} = V_I(t) \tilde{U}(t), \quad \tilde{U}(0) = I
\]

\textbf{Step 2: 计算 $V_I(t)$}

利用海森堡演化求导公式：
\[
\frac{d}{dt}\left( e^{iH_0 t/\hbar} X e^{-iH_0 t/\hbar} \right) = \frac{i}{\hbar} e^{iH_0 t/\hbar} [H_0, X] e^{-iH_0 t/\hbar} = \frac{P}{m}
\]
因此：
\[
e^{iH_0 t/\hbar} X e^{-iH_0 t/\hbar} = X + \frac{P}{m}t
\]
得到相互作用势：
\[
V_I(t) = -\gamma \left( X + \frac{P}{m}t \right)
\]

\textbf{Step 3: 代入验证并确定 $\alpha(t)$}

设 $\tilde{U} = e^{-i\alpha/\hbar} e^{i\beta P} e^{i\eta X}$，其中 $\beta = \frac{\gamma t^2}{2m\hbar}$，$\eta = \frac{\gamma t}{\hbar}$。

计算左式 $i\hbar \frac{d\tilde{U}}{dt}$ 的三项贡献：
\begin{enumerate}[nosep]
    \item 相位项：$\dot{\alpha} \tilde{U}$
    \item $P$ 指数项：$-\hbar\dot{\beta} P \tilde{U} = -\frac{\gamma t}{m} P \tilde{U}$
    \item $X$ 指数项：利用 $e^{i\beta P} X e^{-i\beta P} = X + \beta\hbar = X + \frac{\gamma t^2}{2m}$，得 $-\hbar\dot{\eta}(X + \frac{\gamma t^2}{2m})\tilde{U} = -\gamma(X + \frac{\gamma t^2}{2m})\tilde{U}$
\end{enumerate}

合并得：
\[
i\hbar \frac{d\tilde{U}}{dt} = \left[ \dot{\alpha} - \frac{\gamma t}{m}P - \gamma X - \frac{\gamma^2 t^2}{2m} \right] \tilde{U}
\]

令其等于 $V_I(t)\tilde{U} = (-\gamma X - \frac{\gamma t}{m}P)\tilde{U}$，对比常数项：
\[
\dot{\alpha} = \frac{\gamma^2 t^2}{2m} \quad \Rightarrow \quad \alpha(t) = \frac{\gamma^2 t^3}{6m}
\]
（取积分常数 $\alpha(0) = 0$）

\textbf{Step 4: 薛定谔绘景的 $U(t)$}

由 $U(t) = e^{-iH_0 t/\hbar} \tilde{U}(t)$ 得：
\[
\begin{aligned}
U(t) &= \exp\left( -\frac{iP^2 t}{2m\hbar} \right) \exp\left( -\frac{i\gamma^2 t^3}{6m\hbar} \right) \\
&\quad \times \exp\left( \frac{i\gamma P t^2}{2m\hbar} \right) \exp\left( \frac{i\gamma X t}{\hbar} \right)
\end{aligned}
\]
\end{derivationbox}

\begin{formulabox}[核心结论]
\[
\alpha(t) = \frac{\gamma^2 t^3}{6m}
\]
\[
V_I(t) = -\gamma \left( X + \frac{P}{m}t \right)
\]
\[
\begin{aligned}
U(t) &= e^{-iH_0 t/\hbar} \exp\left( -\frac{i\gamma^2 t^3}{6m\hbar} \right) \\
&\quad \times \exp\left( \frac{i\gamma P t^2}{2m\hbar} \right) \exp\left( \frac{i\gamma X t}{\hbar} \right)
\end{aligned}
\]
\end{formulabox}

\begin{errorbox}[colback=white]{思路短路与注意力分配策略}
\textbf{短路点：}纠结于解的形式如何``天降''，试图从零推导三个指数的乘积结构。

\textbf{策略：}\key{放弃推导，拥抱验证。}在考场上，如果题目给出了具体的解的形式，你的任务不是去``发明''它，而是作为一名质检员，通过求导把未知参数确定下来。

\textbf{注意力分配：}
\begin{itemize}[nosep]
    \item $0\sim20\%$：写出 $V_I(t) = -\gamma(X + \frac{P}{m}t)$
    \item $20\sim80\%$：进行算符穿梭（$e^{i\beta P} X e^{-i\beta P} = X + \beta\hbar$）和系数对齐
    \item $80\sim100\%$：检查 $\alpha(t)$ 的系数。$t^2$ 的积分必须出 $\frac{1}{3}$，再与前面的 $\frac{1}{2}$ 配合，最终得到 $\frac{1}{6}$
\end{itemize}
\end{errorbox}

\begin{insightbox}[举一反三与物理直觉]
\textbf{物理直觉——自由落体的相因子：}

线性势（重力场/匀强电场）中，粒子的量子相位随时间以 $t^3$ 增长。这与经典力学中的作用量 $S = \int L dt$ 直接对应。对于经典轨迹 $x(t) = x_0 + v_0 t + \frac{\gamma t^2}{2m}$，拉格朗日量 $L = \frac{1}{2}m\dot{x}^2 + \gamma x$ 的积分恰好产生 $t^3$ 项。在原子干涉仪（\en{atom interferometry}）中，这一相位是测量重力加速度 $g$ 的核心物理量。

\textbf{拓展延伸：}

\textit{(1) 路径积分视角}：线性势的传播子是可以精确计算的（涉及 Airy 函数），其相位中 $t^3$ 项对应于经典作用量。这是费曼路径积分中最简单的可解析模型之一。

\textit{(2) 定态图像}：含时演化算符的傅里叶变换与定态薛定谔方程的 Airy 函数解（$\text{Ai}(z)$）直接相关。能量本征态在动量空间是平面波，在坐标空间是被线性势``弯曲''的 Airy 函数。

\textit{(3) 与 WKB 的联系}：线性势的 WKB 近似精确给出 Airy 函数的渐近形式，其转折点附近的相位累积由 $t^3$ 型项主导。
\end{insightbox}

\begin{thinkbox}[解题思路模板（绘景变换与算符演化）]
\begin{enumerate}
    \item \textbf{拆分哈密顿量}：明确 $H_0$ 和 $V$，写出相互作用绘景定义
    \item \textbf{计算 $V_I(t)$}：利用 $\frac{d}{dt}(e^{iH_0 t/\hbar} O e^{-iH_0 t/\hbar}) = \frac{i}{\hbar}e^{iH_0 t/\hbar}[H_0, O]e^{-iH_0 t/\hbar}$
    \item \textbf{识别题型}：若题目给出 $\tilde{U}$ 的显式形式，立即启动代入验证
    \item \textbf{算符平移}：熟练使用 $e^{i\beta P} X e^{-i\beta P} = X + \beta\hbar$
    \item \textbf{系数对齐}：分别对比 $X$ 项、$P$ 项和常数（c 数）项
\end{enumerate}
\end{thinkbox}

\end{multicols}

% ============================================
% 补充说明页：线性势场
% ============================================
\newpage
\section{补充说明：线性势场与相互作用绘景的完整细节}

\subsection{题型反射弧标签}

\begin{center}
\begin{tabular}{|c|l|}
\hline
\textbf{层级} & \textbf{识别标签} \\
\hline
一级（领域） & 量子力学，相互作用绘景 \\
二级（方法） & 绘景变换，算符演化，BCH 公式 \\
三级（关键词） & interaction picture, linear potential, Magnus expansion \\
\hline
\end{tabular}
\end{center}

\examnote{看到 $H = H_0 + V$ 并要求相互作用绘景，立即写出 $|\tilde{\psi}\rangle = e^{iH_0 t/\hbar}|\psi\rangle$ 和 $i\hbar \dot{\tilde{U}} = V_I \tilde{U}$。}

\subsection{$V_I(t)$ 的详细推导}

$H_0 = \frac{P^2}{2m}$，$V = -\gamma X$。我们需要计算：
\[
V_I(t) = e^{iH_0 t/\hbar} (-\gamma X) e^{-iH_0 t/\hbar} = -\gamma X_I(t)
\]
其中 $X_I(t)$ 是 $X$ 在相互作用绘景中的演化。对参数 $t$ 求导：
\[
\frac{dX_I}{dt} = \frac{i}{\hbar} e^{iH_0 t/\hbar} [H_0, X] e^{-iH_0 t/\hbar}
\]
计算对易子：
\[
[H_0, X] = \frac{1}{2m}[P^2, X] = \frac{1}{2m}(P[P, X] + [P, X]P) = \frac{1}{2m}(-2i\hbar P) = -\frac{i\hbar P}{m}
\]
代入得：
\[
\frac{dX_I}{dt} = \frac{i}{\hbar} \left( -\frac{i\hbar P}{m} \right) = \frac{P}{m}
\]
注意 $P$ 与 $H_0$ 对易，因此 $P$ 在相互作用绘景中不变。积分并利用初始条件 $X_I(0) = X$：
\[
X_I(t) = X + \frac{P}{m}t
\]
因此：
\[
V_I(t) = -\gamma \left( X + \frac{P}{m}t \right)
\]

\subsection{验证法的完整步骤（考场推荐）}

设 $\tilde{U} = e^{-i\alpha/\hbar} e^{i\beta P} e^{i\eta X}$，其中：
\[
\beta(t) = \frac{\gamma t^2}{2m\hbar}, \quad \eta(t) = \frac{\gamma t}{\hbar}
\]

\textbf{第一步：分别求导}
\begin{align*}
\frac{d\tilde{U}}{dt} &= \left( -\frac{i\dot{\alpha}}{\hbar} \right) \tilde{U} + e^{-i\alpha/\hbar} \left( i\dot{\beta} P e^{i\beta P} \right) e^{i\eta X} + e^{-i\alpha/\hbar} e^{i\beta P} \left( i\dot{\eta} X e^{i\eta X} \right) \\
&= \left( -\frac{i\dot{\alpha}}{\hbar} \right) \tilde{U} + i\dot{\beta} P \tilde{U} + i\dot{\eta} \underbrace{e^{i\beta P} X e^{-i\beta P}}_{= X + \beta\hbar} \tilde{U}
\end{align*}

\textbf{第二步：乘以 $i\hbar$ 并整理}
\begin{align*}
i\hbar \frac{d\tilde{U}}{dt} &= \dot{\alpha} \tilde{U} - \hbar\dot{\beta} P \tilde{U} - \hbar\dot{\eta}(X + \beta\hbar) \tilde{U} \\
&= \left[ \dot{\alpha} - \frac{\gamma t}{m} P - \gamma X - \frac{\gamma^2 t^2}{2m} \right] \tilde{U}
\end{align*}
（其中 $\dot{\beta} = \frac{\gamma t}{m\hbar}$，$\dot{\eta} = \frac{\gamma}{\hbar}$，$\beta\hbar = \frac{\gamma t^2}{2m}$）

\textbf{第三步：与 $V_I \tilde{U}$ 对比}
\[
V_I \tilde{U} = \left( -\gamma X - \frac{\gamma t}{m} P \right) \tilde{U}
\]
$X$ 项和 $P$ 项自动匹配。常数项要求：
\[
\dot{\alpha} - \frac{\gamma^2 t^2}{2m} = 0 \quad \Rightarrow \quad \dot{\alpha} = \frac{\gamma^2 t^2}{2m}
\]
积分得：
\[
\alpha(t) = \int_0^t \frac{\gamma^2 t'^2}{2m} dt' = \frac{\gamma^2 t^3}{6m}
\]

\subsection{Magnus 展开与 BCH 路径（物理本质）}

当 $[V_I(t_1), V_I(t_2)]$ 为 c 数时，Magnus 展开在二阶截断是精确的：
\[
\tilde{U}(t) = \exp\left( \frac{1}{i\hbar} \int_0^t V_I(t') dt' + \frac{1}{2(i\hbar)^2} \int_0^t dt_1 \int_0^{t_1} dt_2 \, [V_I(t_1), V_I(t_2)] \right)
\]

计算对易子：
\[
[V_I(t_1), V_I(t_2)] = \frac{i\hbar\gamma^2}{m}(t_2 - t_1)
\]

第一项（一阶 Magnus）：
\[
\frac{1}{i\hbar} \int_0^t \left( -\gamma X - \frac{\gamma t'}{m} P \right) dt' = \frac{i\gamma X t}{\hbar} + \frac{i\gamma P t^2}{2m\hbar}
\]

第二项（二阶 Magnus）：
\[
\frac{1}{2(i\hbar)^2} \int_0^t dt_1 \int_0^{t_1} dt_2 \, \frac{i\hbar\gamma^2}{m}(t_2 - t_1) = \frac{\gamma^2 t^3}{12m\hbar}
\]

合并后指数部分为：
\[
\frac{i\gamma P t^2}{2m\hbar} + \frac{i\gamma X t}{\hbar} + \frac{\gamma^2 t^3}{12m\hbar}
\]

利用 BCH 公式 $e^{A+B} = e^A e^B e^{-\frac{1}{2}[A,B]}$（$[A,B]$ 为 c 数），令 $A = \frac{i\gamma P t^2}{2m\hbar}$，$B = \frac{i\gamma X t}{\hbar}$，得 $[A,B] = \frac{i\gamma^2 t^3}{2m\hbar}$。整理后即得到题目给出的三指数形式，其中：
\[
\alpha(t) = \frac{\gamma^2 t^3}{6m}
\]

\subsection{考察的完整知识点}

\begin{enumerate}
    \item \textbf{三种绘景}：薛定谔绘景、海森堡绘景、相互作用绘景的定义与转换关系
    \item \textbf{相互作用绘景的运动方程}：$i\hbar \frac{d}{dt}|\tilde{\psi}\rangle = V_I(t)|\tilde{\psi}\rangle$
    \item \textbf{算符的海森堡演化}：$O_I(t) = e^{iH_0 t/\hbar} O e^{-iH_0 t/\hbar}$ 及其微分方程
    \item \textbf{基本对易子}：$[P^2, X] = -2i\hbar P$，$[X, P] = i\hbar$
    \item \textbf{算符平移公式}：$e^{i\beta P} X e^{-i\beta P} = X + \beta\hbar$
    \item \textbf{Magnus 展开与 BCH 公式}：当对易子为 c 数时的精确处理
    \item \textbf{物理图像}：$t^3$ 相位对应经典作用量和原子干涉仪中的重力测量
\end{enumerate}

\subsection{物理直觉：与经典作用量的联系}

经典拉格朗日量 $L = \frac{1}{2}m\dot{x}^2 + \gamma x$。对于初条件 $x(0) = x_0$，$\dot{x}(0) = v_0$ 的轨迹 $x(t) = x_0 + v_0 t + \frac{\gamma t^2}{2m}$，作用量为：
\[
S = \int_0^t L dt' = \frac{1}{2}mv_0^2 t + \gamma x_0 t + \frac{\gamma v_0 t^2}{2m} + \frac{\gamma^2 t^3}{6m}
\]
量子演化算符中的 $e^{-i\gamma^2 t^3/(6m\hbar)}$ 正是对应于经典作用量中与初始条件无关的 $t^3$ 项。这揭示了量子相位与经典轨迹之间的深刻联系。

\subsection{拓展延伸}

\begin{itemize}
    \item \textbf{定态解与 Airy 函数}：$H\psi = E\psi$ 即 $(-\frac{\hbar^2}{2m}\frac{d^2}{dx^2} - \gamma x)\psi = E\psi$，作变量替换 $z = -(\frac{2m\gamma}{\hbar^2})^{1/3}(x + \frac{E}{\gamma})$ 后化为 Airy 方程 $\frac{d^2\psi}{dz^2} = z \psi$，解为 $\psi(x) \propto \text{Ai}(z)$
    \item \textbf{传播子}：费曼路径积分下线性势的传播子可以解析求出，是研究量子-经典对应的重要范例
    \item \textbf{实验联系}：原子干涉仪利用自由落体原子的 $t^3$ 相位差测量局部重力加速度 $g$，精度可达 $10^{-9} g$ 量级
\end{itemize}

\end{document}
